\section{Introduction}
\label{sec:intro}

Two agents are asked to resize the same Kubernetes deployment at the same
time. One doubles the replica count to absorb a load spike; the other adds
three replicas for headroom. Each reads the current count, reasons about what
to do, and writes back a new value. Starting from four replicas, the only
defensible outcomes are $11$ (double, then add) and $14$ (add, then double).
We ran this workload with two concurrent language-model agents and no
coordination. In every one of twenty trials, across two different models, both
agents read~$4$ and the final count was $7$ or $8$: one update was lost.
Neither agent reasoned incorrectly. Each computed the right value from what it
read. The anomaly is a concurrency-control failure, and no amount of prompt
engineering removes it.

\paragraph{Three answers, three costs.}
The database community has seen this problem before, and each classical answer
carries over to agents with a specific cost.

\emph{Make every tool call a transaction.} This is what most agent frameworks
do implicitly. It is fast and wrong: the unit an application wants to be
atomic is the agent's \emph{intent}---read, decide, write---not the individual
call, and the lost update above lives precisely in the gap between calls.

\emph{Make the whole workflow a transaction.} This is correct, but the
workflow includes the agent's reasoning, which is by far its longest part.
Holding locks across model inference serializes exactly the interval one
would most like to overlap. In our measurements, whole-workflow two-phase
locking is barely faster than running the agents one after another
($1.06\times$).

\emph{Let agents run, and repair conflicts afterwards.} Recent systems order
agents by a launch-time priority, apply writes speculatively, and notify an
agent when a premise it read has changed, relying on the agent to repair its
own work~\cite{lyu2026coagent}. This is fast, but it moves the correctness
argument into the model: the protocol is serializable only if a notified agent
correctly decides whether its plan is affected and correctly repairs it. We
measure that assumption directly; it holds in 91\% of notified trials for a
small model and 95\% for a large one, and when it fails on an append-only log
the wrong record cannot be withdrawn.

\paragraph{Chopping, and what breaks it.}
Transaction chopping~\cite{shasha1995chopping} resolves the tension between
the first two answers for classical workloads. A transaction is split into
pieces that execute and commit independently, so that locks are released
between pieces, and a static graph test guarantees that every interleaving of
pieces is equivalent to some serial execution of the original transactions.
For an agent, the natural place to cut is the reasoning gap: commit the read,
think without holding anything, then write.

The classical theorem, however, assumes that the set of transaction programs
is known before execution. An agent violates this assumption in a specific,
consequential way: \emph{it decides what to do next after seeing what it just
read}. A support agent reads an order and only then discovers that the
resolution requires touching a shared gift-card balance. By that point the
read has committed and its result has been returned. The plan that was
certified is no longer the plan being executed.

One can recertify the revised plan, and that is sound. The difficulty is what
to do when recertification fails. Before execution, a failed certification is
always repairable: merge pieces until, in the limit, each workflow is a single
piece. After a prefix has committed, that repair is unavailable, because
merging a committed piece with a future one is not a physical operation. The
question is whether some weaker repair always exists. We show that it does
not, and we characterize exactly when it does.

\paragraph{Contributions.}
\begin{itemize}
\item We formalize \emph{prefix-aware recertification}: admitting a plan
  revision against a batch in which some pieces have already committed and
  only the uncommitted suffix may be regrouped (\S\ref{sec:model},
  \S\ref{sec:theory}).
\item We prove a complete characterization (Theorem~\ref{thm:char}): a
  revision is admissible under future-only repair if and only if sealing
  every workflow's uncommitted remainder into one piece yields an
  SC-cycle-free plan. This reduces admission to a single cheap test and tells
  the runtime when searching finer choppings is pointless.
\item We prove the corresponding impossibility result
  (Corollary~\ref{cor:imp})---waiting for concurrent workflows to finish does
  not rescue an inadmissible revision---and show that admission must be
  serialized (Theorem~\ref{thm:serial}).
\item We build a runtime that enforces the certificate's hypotheses rather
  than trusting them (\S\ref{sec:runtime}): every tool access is mediated and
  checked fail-closed against the footprint of the executing piece, predicate
  reads are ranges so that phantoms are visible to certification, and pieces
  are isolated by locking. We study three locking disciplines inside a piece.
  The third, \ccs, orders lock waits by the same launch-time priority that
  repair-based protocols use for notifications, which makes deadlock
  unreachable while keeping holds short.
\item We check every soundness claim against a brute-force serializability
  oracle (\S\ref{sec:verification}) and evaluate with real language-model
  agents on twelve contended workloads, two models, five baselines and
  three variants of our own, 1{,}682 trials in total (\S\ref{sec:eval}).
\end{itemize}

\paragraph{Summary of results.}
With a small model driving both agents, \ccs is anomaly-free in 98\% of
trials, against 96\% for whole-workflow locking and 93\% for the repair-based
protocol. On the contended workloads it is as fast as the repair-based
protocol ($1.33\times$ versus $1.38\times$ over serial; the difference is
not significant) and significantly faster than whole-workflow locking
($+0.27\times$, faster on eight of eight workloads). Where the structure of
the workload permits a cut, it reaches $2.15\times$ while whole-workflow
locking stays at $0.98\times$. It never deadlocks, and in that run it is the
only protocol whose slowest workload is still faster than serial execution.

We are equally explicit about what does not work. On ten of our twelve
workloads a read-modify-write or scan-then-insert straddles the reasoning gap,
and certification correctly refuses to cut there. Most mid-flight escalations
into shared state are inadmissible: between 4\% and 37\% are admitted in a
controlled study, and 19\% with a large model on live agents. Certified
chopping does not remove the dependence on the agent; it relocates it, from
\emph{repairing a record correctly} to \emph{retrying after a refusal}, and
changes what a failure costs.

We do not claim a new isolation theorem for static choppings, a general ACID
protocol for agents, or the first transactional treatment of agent tool use.
Our contribution is the revision-admission problem and its solution, a runtime
that enforces it, and an evaluation that reports both where the approach wins
and where the right answer is not to chop.
